\documentclass[10pt,a4paper]{amsart}
\usepackage{amsmath,amssymb,amsthm,graphicx}
% For commutative diagrams we use Xy-pic
\usepackage[all]{xy}

% We use 2cell for 2-commutative diagrams.
\xyoption{2cell}
% Theorem environments.
%
\theoremstyle{plain}
\newtheorem{theorem}[subsection]{Theorem}
\newtheorem{proposition}[subsection]{Proposition}
\newtheorem{lemma}[subsection]{Lemma}

\theoremstyle{definition}
\newtheorem{definition}[subsection]{Definition}
\newtheorem{example}[subsection]{Example}
\newtheorem{exercise}[subsection]{Exercise}
\newtheorem{situation}[subsection]{Situation}

\theoremstyle{remark}
\newtheorem{remark}[subsection]{Remark}
\newtheorem{remarks}[subsection]{Remarks}

\numberwithin{equation}{subsection}

% Macros
%
\def\lim{\mathop{\mathrm{lim}}\nolimits}
\def\colim{\mathop{\mathrm{colim}}\nolimits}
\def\Spec{\mathop{\mathrm{Spec}}}
\def\Hom{\mathop{\mathrm{Hom}}\nolimits}
\def\Ext{\mathop{\mathrm{Ext}}\nolimits}
\def\SheafHom{\mathop{\mathcal{H}\!\mathit{om}}\nolimits}
\def\SheafExt{\mathop{\mathcal{E}\!\mathit{xt}}\nolimits}
\def\Sch{\mathit{Sch}}
\def\Mor{\mathop{\mathrm{Mor}}\nolimits}
\def\Ob{\mathop{\mathrm{Ob}}\nolimits}
\def\Sh{\mathop{\mathit{Sh}}\nolimits}
\def\NL{\mathop{N\!L}\nolimits}
\def\CH{\mathop{\mathrm{CH}}\nolimits}
\def\proetale{{pro\text{-}\acute{e}tale}}
\def\etale{{\acute{e}tale}}
\def\QCoh{\mathit{QCoh}}
\def\Ker{\mathop{\mathrm{Ker}}}
\def\Im{\mathop{\mathrm{Im}}}
\def\Coker{\mathop{\mathrm{Coker}}}
\def\Coim{\mathop{\mathrm{Coim}}}

% Boxtimes
%
\DeclareMathSymbol{\boxtimes}{\mathbin}{AMSa}{"02}

%
% Macros for moduli stacks/spaces
%
\def\QCohstack{\mathcal{QC}\!\mathit{oh}}
\def\Cohstack{\mathcal{C}\!\mathit{oh}}
\def\Spacesstack{\mathcal{S}\!\mathit{paces}}
\def\Quotfunctor{\mathrm{Quot}}
\def\Hilbfunctor{\mathrm{Hilb}}
\def\Curvesstack{\mathcal{C}\!\mathit{urves}}
\def\Polarizedstack{\mathcal{P}\!\mathit{olarized}}
\def\Complexesstack{\mathcal{C}\!\mathit{omplexes}}
% \Pic is the operator that assigns to X its picard group, usage \Pic(X)
% \Picardstack_{X/B} denotes the Picard stack of X over B
% \Picardfunctor_{X/B} denotes the Picard functor of X over B
\def\Pic{\mathop{\mathrm{Pic}}\nolimits}
\def\Picardstack{\mathcal{P}\!\mathit{ic}}
\def\Picardfunctor{\mathrm{Pic}}
\def\Deformationcategory{\mathcal{D}\!\mathit{ef}}
\newenvironment{reference}{}{}
\newenvironment{slogan}{}{}
\newenvironment{history}{}{}
\usepackage{cleveref}
\crefname{originalproposition}{original proposition}{original propositions}
\Crefname{originalproposition}{Original proposition}{Original propositions}
\makeatletter
\newlabel{spaces-limits-proposition-approximate}{{07SU}{source}{}{sourcepointer-1}{}}
\newlabel{spaces-limits-proposition-approximate@cref}{{[originalproposition][0][]07SU}{[1][0][]source}{}{}{}}
\makeatother
\crefname{originalsection}{original section}{original sections}
\Crefname{originalsection}{Original section}{Original sections}
\makeatletter
\newlabel{spaces-limits-section-finite-type-quasi-separated}{{0CP5}{source}{}{sourcepointer-2}{}}
\newlabel{spaces-limits-section-finite-type-quasi-separated@cref}{{[originalsection][0][]0CP5}{[1][0][]source}{}{}{}}
\makeatother
\crefname{originalsection}{original section}{original sections}
\Crefname{originalsection}{Original section}{Original sections}
\makeatletter
\newlabel{groupoids-section-groupoids}{{0230}{source}{}{sourcepointer-3}{}}
\newlabel{groupoids-section-groupoids@cref}{{[originalsection][0][]0230}{[1][0][]source}{}{}{}}
\makeatother
\crefname{originalsection}{original section}{original sections}
\Crefname{originalsection}{Original section}{Original sections}
\makeatletter
\newlabel{spaces-limits-remark-finite-type-gives-well-defined-system}{{0CPB}{source}{}{sourcepointer-4}{}}
\newlabel{spaces-limits-remark-finite-type-gives-well-defined-system@cref}{{[originalsection][0][]0CPB}{[1][0][]source}{}{}{}}
\makeatother
\crefname{originallemma}{original lemma}{original lemmas}
\Crefname{originallemma}{Original lemma}{Original lemmas}
\makeatletter
\newlabel{spaces-limits-lemma-morphism-good-diagram-smooth}{{0CPD}{source}{}{sourcepointer-5}{}}
\newlabel{spaces-limits-lemma-morphism-good-diagram-smooth@cref}{{[originallemma][0][]0CPD}{[1][0][]source}{}{}{}}
\makeatother
\crefname{originallemma}{original lemma}{original lemmas}
\Crefname{originallemma}{Original lemma}{Original lemmas}
\makeatletter
\newlabel{spaces-limits-lemma-morphism-good-diagram-flat}{{0CPC}{source}{}{sourcepointer-6}{}}
\newlabel{spaces-limits-lemma-morphism-good-diagram-flat@cref}{{[originallemma][0][]0CPC}{[1][0][]source}{}{}{}}
\makeatother
\crefname{originallemma}{original lemma}{original lemmas}
\Crefname{originallemma}{Original lemma}{Original lemmas}
\makeatletter
\newlabel{spaces-limits-lemma-good-diagram-fibre-product}{{0CPF}{source}{}{sourcepointer-7}{}}
\newlabel{spaces-limits-lemma-good-diagram-fibre-product@cref}{{[originallemma][0][]0CPF}{[1][0][]source}{}{}{}}
\makeatother
\crefname{originallemma}{original lemma}{original lemmas}
\Crefname{originallemma}{Original lemma}{Original lemmas}
\makeatletter
\newlabel{spaces-limits-lemma-morphism-good-diagram}{{0CPA}{source}{}{sourcepointer-8}{}}
\newlabel{spaces-limits-lemma-morphism-good-diagram@cref}{{[originallemma][0][]0CPA}{[1][0][]source}{}{}{}}
\makeatother
\crefname{originaltheorem}{original theorem}{original theorems}
\Crefname{originaltheorem}{Original theorem}{Original theorems}
\makeatletter
\newlabel{algebraic-theorem-smooth-groupoid-gives-algebraic-stack}{{04TK}{source}{}{sourcepointer-9}{}}
\newlabel{algebraic-theorem-smooth-groupoid-gives-algebraic-stack@cref}{{[originaltheorem][0][]04TK}{[1][0][]source}{}{}{}}
\makeatother
\crefname{originallemma}{original lemma}{original lemmas}
\Crefname{originallemma}{Original lemma}{Original lemmas}
\makeatletter
\newlabel{spaces-groupoids-lemma-quotient-stack-functorial}{{046Q}{source}{}{sourcepointer-10}{}}
\newlabel{spaces-groupoids-lemma-quotient-stack-functorial@cref}{{[originallemma][0][]046Q}{[1][0][]source}{}{}{}}
\makeatother
\crefname{originallemma}{original lemma}{original lemmas}
\Crefname{originallemma}{Original lemma}{Original lemmas}
\makeatletter
\newlabel{spaces-groupoids-lemma-criterion-fibre-product}{{04ZN}{source}{}{sourcepointer-11}{}}
\newlabel{spaces-groupoids-lemma-criterion-fibre-product@cref}{{[originallemma][0][]04ZN}{[1][0][]source}{}{}{}}
\makeatother
\crefname{originallemma}{original lemma}{original lemmas}
\Crefname{originallemma}{Original lemma}{Original lemmas}
\makeatletter
\newlabel{spaces-limits-lemma-limit-from-good-diagram}{{0CP9}{source}{}{sourcepointer-12}{}}
\newlabel{spaces-limits-lemma-limit-from-good-diagram@cref}{{[originallemma][0][]0CP9}{[1][0][]source}{}{}{}}
\makeatother
\crefname{originallemma}{original lemma}{original lemmas}
\Crefname{originallemma}{Original lemma}{Original lemmas}
\makeatletter
\newlabel{stacks-properties-lemma-check-immersion-covering}{{0503}{source}{}{sourcepointer-13}{}}
\newlabel{stacks-properties-lemma-check-immersion-covering@cref}{{[originallemma][0][]0503}{[1][0][]source}{}{}{}}
\makeatother
\title{A finite-type quasi-separated algebraic stack embeds closed into a finitely presented stack over a qcqs base}
\author{The Stacks Project Authors}
\date{stacks-project:stacks-limits.tex:local-snapshot-20261011; GNU FDL 1.2 or later}
\begin{document}
\maketitle
\noindent\textit{Editorial note.} The stack morphism is only quasi-separated, not assumed separated. The separated case used in the preparedGabber–Chow theorem is supporting context, not this general statement.\par\medskip
\noindent\textit{Allowed standard prerequisites (editorial).}\begin{itemize}
\item Explicit prerequisite boundary. Algebraic stacks and smooth groupoid presentations, finite-typeqcqs approximation of algebraic spaces, finite-presentation descent of morphisms and diagrams, quotient stacks and local criteria for closed immersions.
\end{itemize}\medskip
\ifdefined\hypertarget\hypertarget{sourcepointer-1}{}\fi
\noindent\textit{Source reference.} Original proposition 07SU: focused-seeds/AGP/sources/stacks-snapshot-20261011/tags.txt, source lines 9914--9914. Stable Stacks tag for an expressly allowed foundational or notational input; the counted construction remains the two concrete core obligations. The retained local core is not replaced by this pointer. See standard\_prerequisites.\par\medskip
\ifdefined\hypertarget\hypertarget{sourcepointer-2}{}\fi
\noindent\textit{Source reference.} Original section 0CP5: focused-seeds/AGP/sources/stacks-snapshot-20261011/tags.txt, source lines 15915--15915. Stable Stacks tag for an expressly allowed foundational or notational input; the counted construction remains the two concrete core obligations. The retained local core is not replaced by this pointer. See standard\_prerequisites.\par\medskip
\ifdefined\hypertarget\hypertarget{sourcepointer-3}{}\fi
\noindent\textit{Source reference.} Original section 0230: focused-seeds/AGP/sources/stacks-snapshot-20261011/tags.txt, source lines 2725--2725. Stable Stacks tag for an expressly allowed foundational or notational input; the counted construction remains the two concrete core obligations. The retained local core is not replaced by this pointer. See standard\_prerequisites.\par\medskip
\ifdefined\hypertarget\hypertarget{sourcepointer-4}{}\fi
\noindent\textit{Source reference.} Original section 0CPB: focused-seeds/AGP/sources/stacks-snapshot-20261011/tags.txt, source lines 15921--15921. Stable Stacks tag for an expressly allowed foundational or notational input; the counted construction remains the two concrete core obligations. The retained local core is not replaced by this pointer. See standard\_prerequisites.\par\medskip
\ifdefined\hypertarget\hypertarget{sourcepointer-5}{}\fi
\noindent\textit{Source reference.} Original lemma 0CPD: focused-seeds/AGP/sources/stacks-snapshot-20261011/tags.txt, source lines 15923--15923. Stable Stacks tag for an expressly allowed foundational or notational input; the counted construction remains the two concrete core obligations. The retained local core is not replaced by this pointer. See standard\_prerequisites.\par\medskip
\ifdefined\hypertarget\hypertarget{sourcepointer-6}{}\fi
\noindent\textit{Source reference.} Original lemma 0CPC: focused-seeds/AGP/sources/stacks-snapshot-20261011/tags.txt, source lines 15922--15922. Stable Stacks tag for an expressly allowed foundational or notational input; the counted construction remains the two concrete core obligations. The retained local core is not replaced by this pointer. See standard\_prerequisites.\par\medskip
\ifdefined\hypertarget\hypertarget{sourcepointer-7}{}\fi
\noindent\textit{Source reference.} Original lemma 0CPF: focused-seeds/AGP/sources/stacks-snapshot-20261011/tags.txt, source lines 15925--15925. Stable Stacks tag for an expressly allowed foundational or notational input; the counted construction remains the two concrete core obligations. The retained local core is not replaced by this pointer. See standard\_prerequisites.\par\medskip
\ifdefined\hypertarget\hypertarget{sourcepointer-8}{}\fi
\noindent\textit{Source reference.} Original lemma 0CPA: focused-seeds/AGP/sources/stacks-snapshot-20261011/tags.txt, source lines 15920--15920. Stable Stacks tag for an expressly allowed foundational or notational input; the counted construction remains the two concrete core obligations. The retained local core is not replaced by this pointer. See standard\_prerequisites.\par\medskip
\ifdefined\hypertarget\hypertarget{sourcepointer-9}{}\fi
\noindent\textit{Source reference.} Original theorem 04TK: focused-seeds/AGP/sources/stacks-snapshot-20261011/tags.txt, source lines 6258--6258. Stable Stacks tag for an expressly allowed foundational or notational input; the counted construction remains the two concrete core obligations. The retained local core is not replaced by this pointer. See standard\_prerequisites.\par\medskip
\ifdefined\hypertarget\hypertarget{sourcepointer-10}{}\fi
\noindent\textit{Source reference.} Original lemma 046Q: focused-seeds/AGP/sources/stacks-snapshot-20261011/tags.txt, source lines 5481--5481. Stable Stacks tag for an expressly allowed foundational or notational input; the counted construction remains the two concrete core obligations. The retained local core is not replaced by this pointer. See standard\_prerequisites.\par\medskip
\ifdefined\hypertarget\hypertarget{sourcepointer-11}{}\fi
\noindent\textit{Source reference.} Original lemma 04ZN: focused-seeds/AGP/sources/stacks-snapshot-20261011/tags.txt, source lines 6471--6471. Stable Stacks tag for an expressly allowed foundational or notational input; the counted construction remains the two concrete core obligations. The retained local core is not replaced by this pointer. See standard\_prerequisites.\par\medskip
\ifdefined\hypertarget\hypertarget{sourcepointer-12}{}\fi
\noindent\textit{Source reference.} Original lemma 0CP9: focused-seeds/AGP/sources/stacks-snapshot-20261011/tags.txt, source lines 15919--15919. Stable Stacks tag for an expressly allowed foundational or notational input; the counted construction remains the two concrete core obligations. The retained local core is not replaced by this pointer. See standard\_prerequisites.\par\medskip
\ifdefined\hypertarget\hypertarget{sourcepointer-13}{}\fi
\noindent\textit{Source reference.} Original lemma 0503: focused-seeds/AGP/sources/stacks-snapshot-20261011/tags.txt, source lines 6486--6486. Stable Stacks tag for an expressly allowed foundational or notational input; the counted construction remains the two concrete core obligations. The retained local core is not replaced by this pointer. See standard\_prerequisites.\par\medskip
\begin{lemma}
\label{lemma-finite-type-closed-in-finite-presentation}
Let $f : \mathcal{X} \to Y$ be a morphism from an algebraic stack
to an algebraic space. Assume:
\begin{enumerate}
\item $f$ is of finite type and quasi-separated,
\item $Y$ is quasi-compact and quasi-separated.
\end{enumerate}
Then there exists a morphism of finite presentation
$f' : \mathcal{X}' \to Y$ and a closed immersion
$\mathcal{X} \to \mathcal{X}'$ of
algebraic stacks over $Y$.
\end{lemma}

\begin{proof}
Write $Y = \lim_{i \in I} Y_i$ as a limit of algebraic spaces
over a directed set $I$ with affine transition morphisms and
with $Y_i$ Noetherian, see
Limits of Spaces, Proposition \ref{spaces-limits-proposition-approximate}.
We will use the material from
Limits of Spaces, Section
\ref{spaces-limits-section-finite-type-quasi-separated}.

\medskip\noindent
Choose a presentation $\mathcal{X} = [U/R]$.
Denote $(U, R, s, t, c, e, i)$ the corresponding groupoid
in algebraic spaces over $Y$. We may and do assume $U$ is affine.
Then $U$, $R$, $R \times_{s, U, t} R$ are quasi-separated
algebraic spaces of finite type over $Y$. We have two morpisms
$s, t : R \to U$, three morphisms
$c : R \times_{s, U, t} R \to R$,
$\text{pr}_1 : R \times_{s, U, t} R \to R$,
$\text{pr}_2 : R \times_{s, U, t} R \to R$,
a morphism $e : U \to R$, and finally a morphism $i : R \to R$.
These morphisms satisfy a list of axioms which are detailed
in Groupoids, Section \ref{groupoids-section-groupoids}.

\medskip\noindent
According to Limits of Spaces, Remark
\ref{spaces-limits-remark-finite-type-gives-well-defined-system}
we can find an $i_0 \in I$ and inverse systems
\begin{enumerate}
\item $(U_i)_{i \geq i_0}$,
\item $(R_i)_{i \geq i_0}$,
\item $(T_i)_{i \geq i_0}$
\end{enumerate}
over $(Y_i)_{i \geq i_0}$ such that
$U = \lim_{i \geq i_0} U_i$,
$R = \lim_{i \geq i_0} R_i$, and
$R \times_{s, U, t} R = \lim_{i \geq i_0} T_i$
and such that there exist morphisms of systems
\begin{enumerate}
\item $(s_i)_{i \geq i_0} : (R_i)_{i \geq i_0} \to (U_i)_{i \geq i_0}$,
\item $(t_i)_{i \geq i_0} : (R_i)_{i \geq i_0} \to (U_i)_{i \geq i_0}$,
\item $(c_i)_{i \geq i_0} : (T_i)_{i \geq i_0} \to (R_i)_{i \geq i_0}$,
\item $(p_i)_{i \geq i_0} : (T_i)_{i \geq i_0} \to (R_i)_{i \geq i_0}$,
\item $(q_i)_{i \geq i_0} : (T_i)_{i \geq i_0} \to (R_i)_{i \geq i_0}$,
\item $(e_i)_{i \geq i_0} : (U_i)_{i \geq i_0} \to (R_i)_{i \geq i_0}$,
\item $(i_i)_{i \geq i_0} : (R_i)_{i \geq i_0} \to (R_i)_{i \geq i_0}$
\end{enumerate}
with
$s = \lim_{i \geq i_0} s_i$,
$t = \lim_{i \geq i_0} t_i$,
$c = \lim_{i \geq i_0} c_i$,
$\text{pr}_1 = \lim_{i \geq i_0} p_i$,
$\text{pr}_2 = \lim_{i \geq i_0} q_i$,
$e = \lim_{i \geq i_0} e_i$, and
$i = \lim_{i \geq i_0} i_i$.
By Limits of Spaces, Lemma
\ref{spaces-limits-lemma-morphism-good-diagram-smooth}
we see that we may assume that $s_i$ and $t_i$ are smooth
(this may require increasing $i_0$).
By Limits of Spaces, Lemma
\ref{spaces-limits-lemma-morphism-good-diagram-flat}
we may assume that the maps
$R \to U \times_{U_i, s_i} R_i$ given by $s$ and $R \to R_i$ and
$R \to U \times_{U_i, t_i} R_i$ given by $t$ and $R \to R_i$
are isomorphisms for all $i \geq i_0$.
By Limits of Spaces, Lemma \ref{spaces-limits-lemma-good-diagram-fibre-product}
we see that we may assume that the diagrams
$$
\xymatrix{
T_i \ar[r]_{q_i} \ar[d]_{p_i} & R_i \ar[d]^{t_i} \\
R_i \ar[r]^{s_i} & U_i
}
$$
are cartesian. The uniqueness of
Limits of Spaces, Lemma \ref{spaces-limits-lemma-morphism-good-diagram}
then guarantees that for a sufficiently large $i$
the relations between the morphisms $s, t, c, e, i$ mentioned above
are satisfied by $s_i, t_i, c_i, e_i, i_i$. Fix such an $i$.

\medskip\noindent
It follows that $(U_i, R_i, s_i, t_i, c_i, e_i, i_i)$
is a smooth groupoid in algebraic spaces over $Y_i$.
Hence $\mathcal{X}_i = [U_i/R_i]$ is an algebraic stack
(Algebraic Stacks, Theorem
\ref{algebraic-theorem-smooth-groupoid-gives-algebraic-stack}).
The morphism of groupoids
$$
(U, R, s, t, c, e, i) \to (U_i, R_i, s_i, t_i, c_i, e_i, i_i)
$$
over $Y \to Y_i$ determines a commutative diagram
$$
\xymatrix{
\mathcal{X} \ar[d] \ar[r] & \mathcal{X}_i \ar[d] \\
Y \ar[r] & Y_i
}
$$
(Groupoids in Spaces, Lemma
\ref{spaces-groupoids-lemma-quotient-stack-functorial}).
We claim that the morphism $\mathcal{X} \to Y \times_{Y_i} \mathcal{X}_i$
is a closed immersion. The claim finishes the proof because
the algebraic stack $\mathcal{X}_i \to Y_i$ is of finite presentation
by construction. To prove the claim, note that the left diagram
$$
\xymatrix{
U \ar[d] \ar[r] & U_i \ar[d] \\
\mathcal{X} \ar[r] & \mathcal{X}_i
}
\quad\quad
\xymatrix{
U \ar[d] \ar[r] & Y \times_{Y_i} U_i \ar[d] \\
\mathcal{X} \ar[r] & Y \times_{Y_i} \mathcal{X}_i
}
$$
is cartesian by Groupoids in Spaces, Lemma
\ref{spaces-groupoids-lemma-criterion-fibre-product}
and the results mentioned above.
Hence the right commutative diagram is cartesian too.
Then the desired result follows from the fact that
$U \to Y \times_{Y_i} U_i$ is a closed immersion
by construction of the inverse system $(U_i)$ in
Limits of Spaces, Lemma \ref{spaces-limits-lemma-limit-from-good-diagram},
the fact that $Y \times_{Y_i} U_i \to Y \times_{Y_i} \mathcal{X}_i$
is smooth and surjective, and Properties of Stacks, Lemma
\ref{stacks-properties-lemma-check-immersion-covering}.
\end{proof}

\bibliographystyle{plain}
\bibliography{my}
\end{document}
